% Tài liệu tham khảo được in bằng biblatex (APA, biber) từ tailieu.bib và references/chN.bib.
% Theo mẫu của Trường, tài liệu được chia thành hai nhóm: tài liệu tiếng Việt (gắn
% keywords = {tiengviet} trong tệp .bib) và tài liệu tiếng nước ngoài (mọi mục còn lại).
\renewcommand\bibname{DANH MỤC TÀI LIỆU THAM KHẢO}
\chapter*{\centering \MakeUppercase{\bibname}}
\phantomsection
\addcontentsline{toc}{chapter}{\bibname}

\noindent\textbf{Tiếng Việt}
\printbibliography[heading=none, keyword=tiengviet]

\vspace{0.5cm}
\noindent\textbf{Tiếng nước ngoài}
\printbibliography[heading=none, notkeyword=tiengviet]
